% ---------------------------------------------------------------------
\section{Discussion des cas d'échec}
\label{sec:discussion}
% ---------------------------------------------------------------------

\subsection{\texorpdfstring{Pourquoi le MP3 à \qty{64}{\kilo\bit\per\second} gêne moins qu'un bruit blanc à \qty{10}{\decibel} (Q6)}{Pourquoi le MP3 à 64 kbit/s gêne moins qu'un bruit blanc à 10 dB (Q6)}}
\label{sec:q6}

\aremplir{Constat mesuré : \TBD{} = x~\% en MP3 \qty{64}{\kilo\bit\per\second} contre y~\% sous
bruit blanc à \qty{10}{\decibel} ($\Tq$, $\sthr$, nombre de requêtes), et $p$ = x contre y
(\cref{tab:pic-paire}).} Deux dégradations de « niveau » comparable ont donc des effets très
différents. L'explication tient à la manière dont chacune répartit son bruit dans le plan
temps-fréquence, comparée à la définition (4) des pics.

\paragraph{Le codage perceptuel épargne les pics.} Un codeur MP3 découpe le signal en
sous-bandes, estime pour chaque trame un \emph{seuil de masquage}, c'est-à-dire le niveau
en dessous duquel l'oreille ne perçoit rien en présence des composantes fortes voisines, et
alloue ses bits de façon à maintenir le bruit de quantification sous ce
seuil~\cite{painter2000}. Les composantes situées sous le seuil sont grossièrement
quantifiées ou supprimées. Or ce seuil est construit à partir des \emph{masqueurs},
c'est-à-dire des composantes les plus fortes de leur voisinage temps-fréquence ; il se
situe typiquement de \qtyrange{10}{20}{\decibel} sous un masqueur tonal. Un pic au sens de
(4) est, par définition, le maximum de son voisinage $V$ : c'est précisément un masqueur.
Le codeur le conserve, et le bruit qu'il ajoute autour reste nettement plus bas que lui.
Dans les termes de la \cref{sec:maxima}, l'avance $\mu$ du pic sur ses voisins n'est que
peu entamée, et il reste en place. Le codage supprime surtout les composantes faibles,
qui ne sont pas des pics : il simplifie le spectrogramme sans toucher à sa constellation.
Enfin, le bruit de quantification suit l'enveloppe spectrale du signal : le rapport signal
à bruit \emph{local} est à peu près le même dans toutes les bandes.

\paragraph{Le bruit blanc noie les pics des bandes hautes.} Un bruit blanc de variance
$\sigma_b^2$ ajoute en chaque point du spectrogramme une contribution de puissance moyenne
\[
  \E\bigl|B[m,k]\bigr|^2 = \sigma_b^2 \sum_{n=0}^{N-1} w[n]^2 ,
\]
\emph{la même dans tous les canaux}. Le spectre musical, lui, décroît avec la fréquence :
l'essentiel de l'énergie est sous \qty{1}{\kilo\hertz}. Le \RSB{} de \qty{10}{\decibel} est
un rapport \emph{global}, dominé par les basses fréquences ; le rapport local dans le canal
$k$ vaut
\[
  \RSB_{\text{local}}(k) = \RSB + 10\log_{10} \frac{S_x(k)}{\overline{S_x}} ,
\]
où $\overline{S_x}$ est la moyenne du spectre sur tous les canaux. Dans une bande où le
spectre musical est \qty{20}{\decibel} sous sa moyenne, ce qui est courant entre 2 et
\qty{4}{\kilo\hertz}, le rapport local tombe à \qty{-10}{\decibel} : le bruit y domine, et
ses propres fluctuations créent des maxima locaux qui remplacent ceux de la musique. Le
quota par bande aggrave l'effet : il impose de prendre le même nombre de pics dans les
bandes hautes que dans les basses, et ce sont alors des pics de bruit, différents à chaque
requête. Une bonne partie des pics des deux ou trois bandes supérieures est perdue, donc
$p$ chute, et le score, qui varie en $p^2$, chute encore plus.

\paragraph{En résumé.} Le MP3 est conçu pour préserver ce qui domine localement, et c'est
exactement ce que retiennent les pics. Le bruit blanc ne tient aucun compte du signal, et
son effet est maximal là où la musique est faible. \aremplir{Si la mesure de $p$ par bande
est disponible, l'ajouter ici : $p$ à peu près constant selon la bande en MP3, décroissant
vers les hautes fréquences sous bruit blanc.} À \qty{32}{\kilo\bit\per\second}, le codeur
manque de bits : le bruit de quantification dépasse le seuil de masquage, les composantes
moyennes disparaissent par blocs et certains pics secondaires avec elles, ce qui explique
la baisse du \TBD{} à bas débit (\cref{fig:tbd-mp3}).

\subsection{Genres difficiles}
\label{sec:genres}

\aremplir{Constat mesuré : \TBD{} par genre (\cref{fig:genres}), genres les plus et les
moins bien reconnus, avec $\Tq$, $\sthr$ et la dégradation.}

La théorie permet d'anticiper quels genres posent problème, et pourquoi.

\begin{itemize}
  \item \emph{Musiques tonales et tenues (Folk, Instrumental, une partie d'International).}
        Les notes sont stables, leurs harmoniques dépassent nettement le fond : les pics ont
        une forte avance $\mu$ et survivent bien au bruit et au filtrage. On attend un bon
        \TBD{}. En contrepartie, les notes tenues produisent des clés répétées sur plusieurs
        trames, donc des votes fortuits cohérents (\cref{sec:decision}) : les scores
        négatifs peuvent y être plus élevés.
  \item \emph{Musiques à spectre dense ou bruité (Rock saturé, Experimental).} Une guitare
        saturée ou une texture bruitée remplit le spectre de façon presque uniforme : la
        platitude spectrale est élevée, les maxima locaux dépassent à peine leurs voisins
        ($\mu$ faible), et le moindre bruit les déplace. La condition (5) et son équivalent
        pour le bruit sont rarement remplies : $p$ est faible.
  \item \emph{Musiques dominées par les percussions et la voix parlée (Hip-Hop).} L'énergie
        est concentrée dans des transitoires brefs et dans les basses ; une voix parlée a
        une hauteur qui glisse en permanence. Les pics sont peu nombreux dans les bandes
        moyennes et hautes, ou instables en fréquence, et un décalage non aligné sur la
        grille (\cref{sec:proprietes}) déplace facilement un transitoire d'une trame.
  \item \emph{Musiques en boucles (Electronic, Hip-Hop).} Une boucle répétée produit les
        mêmes clés à chaque répétition : l'histogramme des $\delta$ du bon morceau présente
        plusieurs pics, ce qui ne gêne pas l'identification mais rend la position estimée
        ambiguë. Des sons de synthèse ou des échantillons de batterie partagés entre
        morceaux peuvent en revanche créer des coïncidences cohérentes avec de mauvais
        morceaux.
\end{itemize}
\aremplir{Confronter ces hypothèses à la mesure : les genres prévus comme difficiles le
sont-ils ? Exemple d'échec analysé : spectrogramme et constellation d'une requête ratée.}

\begin{figure}[t]
  \centering
  \inclurefig[0.8\linewidth]{tbd_par_genre}{TBD par genre}
  \caption{\TBD{} par genre pour \aremplir{dégradation, $\Tq$, $\sthr$, nombre de requêtes
  par genre}.}
  \label{fig:genres}
\end{figure}

\subsection{Réverbération et enregistrements réels}
\label{sec:echecs-reverb}

La réverbération sort du cadre de la propriété (2) dès que sa réponse dépasse la durée de
la fenêtre (\cref{sec:limites}). Elle agit comme un bruit additif fabriqué par le signal
lui-même : on attend une dégradation modérée à $T_{60} = \qty{0.3}{\second}$ et nette à
$T_{60} = \qty{1}{\second}$, plus marquée pour les morceaux rapides que pour les morceaux
lents. \aremplir{Résultats de la réverbération et des enregistrements en salle à 1 et
4~m : \TBD{} mesurés, et comparaison avec les dégradations simulées.}

\subsection{Limites de la méthode}

La méthode n'est invariante ni au changement de tempo (les écarts $\dm$ changent) ni au
changement de hauteur (les canaux changent) : une version accélérée, transposée ou jouée
en direct n'est pas reconnue. Elle identifie un \emph{enregistrement}, pas une
\emph{œuvre}.
